\chapter{Ulnar Nerve Release}

\section*{What Is This Surgery?}
Ulnar nerve release, commonly referred to as cubital tunnel decompression, is a surgical intervention aimed at relieving pressure on the ulnar nerve as it traverses the cubital tunnel at the elbow. This condition, known as cubital tunnel syndrome, manifests as numbness and tingling in the ring and little fingers, pain along the medial forearm, and potential weakness or atrophy of the intrinsic muscles of the hand. The surgical procedure is clinically significant as it restores normal nerve function by either enlarging the cubital tunnel or relocating the nerve to a less compressive environment, thus enhancing blood supply and nerve conduction. Ulnar nerve release is particularly indicated for patients experiencing persistent and severe symptoms that have not responded adequately to conservative management strategies.

\section*{Alternative Names}
\begin{itemize}
  \item Cubital tunnel decompression
  \item Ulnar nerve decompression at the elbow
  \item Anterior transposition of the ulnar nerve
  \item Subcutaneous ulnar nerve transposition
  \item Submuscular ulnar nerve transposition
  \item Ulnar neurolysis
\end{itemize}

\section*{Relevant Surgical Anatomy}
The ulnar nerve is a major peripheral nerve that arises from the medial cord of the brachial plexus and descends along the medial aspect of the arm. As it approaches the elbow, it passes posterior to the medial epicondyle of the humerus, entering the cubital tunnel, which is bordered by the medial epicondyle, the olecranon process of the ulna, and the cubital tunnel retinaculum (Osborne's ligament). 

Key anatomical features relevant to this procedure include:

\begin{itemize}
  \item **Arterial Supply:** The ulnar nerve receives blood supply from the ulnar artery, which branches from the brachial artery.
  \item **Venous Drainage:** Venous drainage is primarily via the accompanying veins of the ulnar artery.
  \item **Nerve Relations:** The ulnar nerve is accompanied by the medial antebrachial cutaneous nerve and is in close proximity to the brachial artery and median nerve.
  \item **Lymphatic Drainage:** Lymphatic drainage occurs through the cubital lymph nodes located in the medial aspect of the elbow.
  \item **Surgical Landmarks:** Key landmarks include McBurney's point for abdominal surgeries, but for ulnar nerve release, the medial epicondyle serves as a critical reference point.
\end{itemize}

Understanding these anatomical relationships is crucial for minimizing complications during surgical intervention.

\section*{Surgical Principle \& Rationale}
The primary surgical principle of ulnar nerve release is to alleviate mechanical compression of the ulnar nerve within the cubital tunnel. This compression can be exacerbated by factors such as elbow flexion, direct trauma, or anatomical variations like osteophytes. The operative goal is to decompress the nerve, thereby enhancing its vascularity and facilitating axonal transport.

The two main surgical strategies employed are simple decompression and anterior transposition. Simple decompression involves incising the cubital tunnel retinaculum and any constricting fascial bands to enlarge the tunnel. In contrast, anterior transposition relocates the ulnar nerve from its posterior position to an anterior location, which can be subcutaneous, intramuscular, or submuscular. This relocation prevents the nerve from being stretched or compressed during elbow flexion, particularly in cases where the nerve is prone to subluxation. Both techniques aim to eliminate the source of compression, allowing the nerve to glide freely and facilitating recovery of function.

\section*{History \& Surgical Evolution}
The understanding of ulnar nerve compression at the elbow has evolved significantly since the late 19th century, when early descriptions of symptoms associated with trauma to the "funny bone" region were first documented. In 1908, Sir Harold Stiles performed one of the earliest recorded ulnar nerve transpositions, marking a pivotal moment in the surgical management of this condition.

In 1942, Learmonth introduced the submuscular transposition technique, which involved detaching the flexor-pronator muscle origin, placing the nerve beneath it, and reattaching the muscle to provide a well-vascularized bed for the nerve. Over the decades, various techniques have been refined, including subcutaneous and intramuscular transposition methods, as well as the increasing popularity of simple decompression. The evolving understanding of the dynamic nature of ulnar nerve compression during elbow flexion has significantly influenced modern surgical approaches.

\section*{Clinical Indications}
Ulnar nerve release is indicated in the following clinical scenarios:

\begin{itemize}
  \item Persistent ulnar neuropathy symptoms, such as numbness, tingling, or pain in the ring and little fingers, that have failed to improve with conservative management.
  \item Progressive motor weakness or muscle atrophy in the ulnar-innervated intrinsic hand muscles, indicating significant nerve compromise.
  \item Positive electrodiagnostic studies confirming ulnar nerve compression at the elbow and quantifying its severity.
  \item Documented subluxation or dislocation of the ulnar nerve over the medial epicondyle with elbow flexion, leading to repetitive microtrauma.
  \item Presence of a space-occupying lesion causing direct compression of the ulnar nerve.
  \item Acute nerve injury with evidence of ongoing compression or entrapment.
  \item Symptoms significantly impairing daily activities, work, or sleep quality.
\end{itemize}

\section*{Clinical Positioning}
Ulnar nerve release is typically considered a secondary procedure, performed after conservative management has failed to alleviate symptoms or prevent progression of nerve dysfunction. Conservative treatments may include activity modification, elbow splinting, NSAIDs, and physical therapy. However, in cases of severe or rapidly progressive ulnar neuropathy, surgery may be warranted as a primary intervention to prevent irreversible nerve damage. The decision to proceed with surgery is based on symptom severity, duration, and findings from electrodiagnostic studies.

\section*{Surgical Technique \& Procedure}
Ulnar nerve release is generally performed under regional anesthesia with intravenous sedation or general anesthesia. The patient is positioned supine with the arm abducted and externally rotated, and the elbow flexed to approximately 70-90 degrees for optimal exposure.

The surgical procedure involves the following key steps:

\begin{itemize}
  \item A curvilinear incision is made over the medial epicondyle, extending proximally and distally to adequately expose the ulnar nerve.
  \item Careful dissection is performed to identify the ulnar nerve both proximal and distal to the cubital tunnel, ensuring its protection throughout the procedure.
  \item The cubital tunnel retinaculum and any constricting fascial bands are released to enlarge the tunnel.
  \item Neurolysis may be performed to free the nerve from surrounding scar tissue or adhesions.
  \item The decision is made between simple decompression and anterior transposition:
    \begin{itemize}
      \item For simple decompression, the nerve is left in its anatomical bed after all constricting structures are released.
      \item For anterior transposition, the nerve is mobilized and moved to an anterior position, which can be subcutaneous, intramuscular, or submuscular.
    \end{itemize}
  \item Adequate hemostasis is ensured to prevent hematoma formation.
  \item Closure of the subcutaneous tissues and skin is performed with sutures.
  \item A soft dressing is applied, often with a posterior splint or brace to protect the elbow and limit flexion.
\end{itemize}

\section*{Preoperative Workup \& Preparation}
Thorough preoperative preparation is essential for a successful outcome. This typically includes:

\begin{itemize}
  \item A detailed medical history and physical examination to assess the extent of ulnar neuropathy and identify comorbidities.
  \item Electrodiagnostic studies to confirm the diagnosis and assess the severity of nerve damage.
  \item Imaging studies, such as plain X-rays and possibly MRI, to rule out bony abnormalities or soft tissue masses.
  \item Routine preoperative blood tests, including a complete blood count and coagulation profile.
  \item Cardiac and anesthesia clearance for patients with significant medical conditions.
  \item Instructions to be NPO for 6-8 hours prior to surgery.
  \item Adjustment or temporary discontinuation of anticoagulant medications as per surgeon's instructions.
  \item Administration of prophylactic antibiotics intravenously prior to incision.
  \item Obtaining comprehensive informed consent, ensuring the patient understands the procedure, risks, benefits, and alternatives.
\end{itemize}

\section*{Contraindications \& Precautions}
\textbf{Absolute Contraindications:}
\begin{itemize}
  \item Active infection at the surgical site or systemic sepsis.
  \item Uncontrolled coagulopathy or bleeding disorder.
  \item Severe systemic illness or unstable medical conditions.
  \item Patient refusal or inability to provide informed consent.
\end{itemize}

\textbf{Relative Contraindications and Precautions:}
\begin{itemize}
  \item Mild, intermittent symptoms well-controlled with conservative measures.
  \item Poor patient compliance with postoperative rehabilitation protocols.
  \item Severe, long-standing ulnar neuropathy with irreversible nerve damage.
  \item Significant comorbidities such as uncontrolled diabetes or severe peripheral vascular disease.
  \item Smoking, associated with impaired wound healing.
  \item Previous surgery in the same area, complicating dissection due to scar tissue.
\end{itemize}

\section*{Complications \& Risks}
Ulnar nerve release carries potential risks, categorized as general surgical risks and procedure-specific risks.

\textbf{General Surgical Risks:}
\begin{itemize}
  \item Bleeding: Hematoma formation at the surgical site.
  \item Infection: Superficial or deep wound infection.
  \item Anesthesia complications: Nausea, allergic reactions, respiratory or cardiac complications.
  \item Scarring: Formation of a noticeable scar at the incision site.
  \item Pain: Persistent or chronic pain at the surgical site.
\end{itemize}

\textbf{Procedure-Specific Risks:}
\begin{itemize}
  \item Persistent or recurrent symptoms due to incomplete decompression.
  \item Damage to the ulnar nerve during dissection.
  \item Injury to the medial antebrachial cutaneous nerve.
  \item Elbow stiffness postoperatively.
  \item Complex Regional Pain Syndrome (CRPS).
  \item Failure of transposition leading to recurrent symptoms.
  \item Neuroma formation at the site of nerve injury.
  \item Vascular injury, though rare.
\end{itemize}

\section*{Postoperative Recovery Timeline}
Following ulnar nerve release, patients are monitored in the Post-Anesthesia Care Unit (PACU) for vital signs, pain management, and neurovascular checks. The arm is typically elevated to reduce swelling, and pain medication is administered as needed. Most patients are discharged home the same day or the following morning.

During the first 24-48 hours, patients should keep the incision clean and dry, continuing prescribed pain medication. A sling or soft dressing is often used for comfort and protection. Suture removal usually occurs around 10-14 days postoperatively. Gentle range of motion exercises for the elbow and hand may begin shortly after surgery, particularly if a simple decompression was performed. For transposition procedures, a period of immobilization in a splint for 2-3 weeks may be recommended, followed by a gradual increase in activity and physical therapy. Full resolution of numbness and tingling can take several months, as nerve regeneration is a slow process.

\section*{Surgical Findings \& Pathology}
During the surgical procedure, the surgeon documents the condition of the ulnar nerve and any surrounding structures. The pathology report is typically not relevant unless a mass or abnormal tissue was excised. The primary focus is on the successful release of the nerve from any constricting structures, allowing for improved function postoperatively.

\section*{Expected Postoperative Course}
A normal postoperative course following ulnar nerve release includes a gradual reduction in symptoms such as pain, numbness, and tingling. Patients typically experience a timeline of healing that includes initial discomfort, followed by a return of sensation and strength in the affected hand muscles. While complete recovery may take several months, most patients report significant improvement in their symptoms and functional abilities.

\section*{Postoperative Care \& Follow-up}
Postoperative care includes regular follow-up visits to monitor recovery and assess for any complications. Key components of postoperative care include:

\begin{itemize}
  \item Post-operative visit for wound check and suture removal, typically 10-14 days after surgery.
  \item Follow-up visits at 4-6 weeks and 3-6 months to assess symptom resolution and functional recovery.
  \item Physical therapy may be prescribed to aid in regaining full elbow range of motion and strength.
  \item Gradual resumption of light activities, progressing to more strenuous tasks as tolerated.
  \item Patients should monitor for warning signs such as increasing pain, redness, swelling, fever, or new/worsening numbness or weakness, and contact their surgeon if these occur.
\end{itemize}

\section*{Epidemiology}
Cubital tunnel syndrome is the second most common peripheral nerve entrapment neuropathy, with an estimated incidence of 20 to 30 cases per 100,000 person-years. It affects men more frequently than women, with prevalence increasing with age, particularly in individuals over 50 years old. Common risk factors include repetitive elbow flexion, prolonged leaning on the elbows, direct trauma, and systemic conditions such as diabetes mellitus. Untreated cubital tunnel syndrome can lead to significant morbidity, including chronic pain and functional impairment.

\section*{Outcomes \& Prognosis}
The outcomes of ulnar nerve release are generally favorable, with success rates ranging from 70\% to 90\% in achieving significant symptom relief and functional improvement. Both simple decompression and anterior transposition techniques have demonstrated comparable efficacy, though the choice of procedure often depends on individual patient factors. Patients with shorter symptom duration and less severe preoperative deficits tend to have better prognoses. While complete resolution of severe motor weakness may not always occur, most patients experience substantial improvement in pain and function. Recurrence rates are relatively low, typically ranging from 5\% to 10\%, with factors such as diabetes and smoking identified as potential negative prognostic indicators.

\section*{References}
\begin{enumerate}
  \item MedlinePlus. Cubital tunnel syndrome. U.S. National Library of Medicine; 2024.
  \item Canale ST, Beaty JH. Campbell's Operative Orthopaedics. 14th ed. Elsevier; 2021.
  \item American Academy of Orthopaedic Surgeons (AAOS). Management of Cubital Tunnel Syndrome. Clinical Practice Guideline; 2016.
  \item Netter FH. Atlas of Human Anatomy. 8th ed. Elsevier; 2023.
\end{enumerate}

\clearpage